\chapter*{}
\addcontentsline{toc}{chapter}{Resumen y Palabras Clave}
% Resumen en español
\Large{\textsc{\textcolor{nar}{Resumen}}} \par
\vspace{0.2cm}
\normalsize{%
Este Trabajo de Fin de Máster presenta el desarrollo de una plataforma de comunicación con
periféricos sobre el MPSoC Zynq UltraScale+ ZCU102, en el marco del proyecto LINCE, una
iniciativa del PERTE Aeroespacial liderada por Indra para el desarrollo de microsatélites de
órbita baja.

El trabajo abarca cinco áreas principales.

En primer lugar, el diseño e implementación en VHDL de un transceptor serie altamente configurable para 
la lógica programable del MPSoC.

En segundo lugar, el desarrollo de un driver serie en C sobre el sistema operativo en tiempo real RTEMS 
en el sistema de procesamiento del MPSoC, que utiliza el transceptor serie desarrollado y ofrece una 
API pública que abstrae completamente del hardware. 
En tercer lugar, el diseño y fabricación de tres placas de circuito impreso de ampliación de la ZCU102
que se conectan a ella para expandir sus capacidades y permiten ampliar la ventana de
desarrollo. La primera placa ofrece soporte hardware para 7 líneas RS422 y 7 líneas RS485,
las cuales se pueden conectar en buses compartidos configurables por hardware. La segunda
placa ofrece soporte para 3 líneas RS422 o RS485 seleccionable, dos líneas CAN, cuatro
líneas PWM y cuatro líneas analógicas. La tercera placa ofrece soporte para cinco líneas
RS422 o RS485 seleccionable, tres pares de líneas de potencia para controlar motores por
PWM y dos líneas SpaceWire.

En cuarto lugar, el estudio del transporte de datos entre el sistema de procesamiento y la lógica
programable, que resultó ser el factor determinante del rendimiento del conjunto. Se implementaron
y midieron tres arquitecturas completas sobre el mismo transceptor: una basada en AXI GPIO, otra
con catorce controladores de acceso directo a memoria (DMA) y una tercera con un único controlador
multicanal (MCDMA) gobernado por un puente diseñado en VHDL. La comparativa demuestra que la
primera interrumpe al procesador una vez por byte y se queda en el 30\% del techo de la línea,
mientras que la tercera alcanza el 99\% con cuatrocientas veces menos interrupciones.

En quinto lugar, la validación de la plataforma mediante el desarrollo de herramientas de software
y hardware adicionales, que permitieron verificar el correcto funcionamiento del driver serie
y de las distintas funcionalidades de las placas de ampliación.

El conjunto del trabajo constituye una plataforma funcional y validada que el equipo de
Sener, Indra y el laboratorio B105 pueden utilizar como base para el subsistema de
comunicaciones del ordenador de a bordo del satélite LINCE en el futuro.
}

\vspace{0.9cm}

% Summary (resumen en inglés)
\Large{\textsc{\textcolor{nar}{Summary}}} \par
\vspace{0.2cm}
\normalsize{%
This Master's Thesis presents the development of a peripheral communication platform on the Zynq
UltraScale+ ZCU102 MPSoC, within the framework of the LINCE project, an initiative of the Spanish
Aerospace PERTE led by Indra for the development of low Earth orbit microsatellites.

The work covers five main areas.

First, the design and implementation in VHDL of a highly configurable serial transceiver for the
programmable logic of the MPSoC.

Second, the development of a serial driver in C on the RTEMS real-time operating system running on
the processing system of the MPSoC, which uses the serial transceiver developed and provides a
public API that fully abstracts the hardware.

Third, the design and manufacture of three expansion printed circuit boards for the ZCU102, which
connect to it in order to extend its capabilities and widen the development window. The first board
provides hardware support for 7 RS422 lines and 7 RS485 lines, which can be connected in shared
buses configurable by hardware. The second board provides support for 3 selectable RS422 or RS485
lines, two CAN lines, four PWM lines and four analogue lines. The third board provides support for
five selectable RS422 or RS485 lines, three pairs of power lines for PWM motor control and two
SpaceWire lines.

Fourth, the study of data transport between the processing system and the programmable logic, which
proved to be the determining factor in the overall performance of the system. Three complete
architectures were implemented and measured on the same transceiver: one based on AXI GPIO, another
with fourteen independent direct memory access (DMA) controllers, and a third one with a single
multichannel controller (MCDMA) governed by a bridge designed in VHDL. The comparison shows that
the first one interrupts the processor once per byte and reaches only 30\% of the physical limit of
the line, whereas the third one reaches 99\% with four hundred times fewer interrupts.

Fifth, the validation of the platform through the development of additional software and hardware
tools, which made it possible to verify the correct operation of the serial driver and of the
different features of the expansion boards.

Taken as a whole, the work constitutes a functional and validated platform that the teams at Sener,
Indra and the B105 laboratory can use in the future as a basis for the communications subsystem of
the on-board computer of the LINCE satellite.
}
\vspace{0.9cm}

% Palabras clave
\Large{\textsc{\textcolor{nar}{Palabras Clave}}} \par \par
\vspace{0.2cm}
\normalsize{MPSoC, FPGA, VHDL, RTEMS, RS422, RS485, CAN, SpaceWire, driver serie, PCB,
AXI, DMA, MCDMA, AXI-Stream, sistemas empotrados, tiempo real, satélite, ordenador de a bordo.} \par
\vspace{0.8cm}

\Large{\textsc{\textcolor{nar}{Keywords}}} \par
\vspace{0.2cm}
\normalsize{MPSoC, FPGA, VHDL, RTEMS, RS422, RS485, CAN, SpaceWire, serial driver, PCB,
AXI, DMA, MCDMA, AXI-Stream, embedded systems, real time, satellite, on-board computer.} \par
